\chapter{项目交接与新病例落地建议}

\section{接手项目的推荐阅读顺序}

\begin{enumerate}
  \item 先读 \projectpath{README.md} 和本报告，建立输入—规划—实体化—验证的整体认识；
  \item 阅读 \projectpath{config.py} 的数据类和 \texttt{CaseConfig.from\_json()}，理解配置
  合并、安全门和模式互斥；
  \item 以 \projectpath{generation_process.py} 为索引，逐步定位七阶段入口；
  \item 重点读 \projectpath{tooth_section_anchors.py}、\projectpath{point_linking.py} 和
  \projectpath{blender/guide_modeling.py}，它们决定锚点、中心线和最终布尔顺序；
  \item 最后读 \projectpath{guide_validation.py} 和测试，理解项目真正承诺检查什么。
\end{enumerate}

\section{新病例的建议流程}

\begin{enumerate}
  \item 将导板、牙列、双导管装配体放入规范病例目录，确认毫米单位和共同坐标系；
  \item 在 YAML 中填写上下颌、FDI 顺序、现存/缺失/排除牙、观察窗和人工确认方向；
  \item 选择 input 或 generated 导管模式，判断是否需要断裂跨接、末端 U 梁或远中公共节点；
  \item 在 YAML 中为每个主梁和 Y 梁端点指定牙位站位、侧别和射线角度；
  \item 先运行 \texttt{process}，检查牙位、T 点、射线、A 点、Q/P 和中心线；
  \item 查看牙位映射、观察窗、锚点和梁架过程图，只修正语义或明确参数，不用无限放大容差；
  \item 在专用输出目录执行完整生成，人工检查四视图和导孔/窗口；
  \item 更新 review 状态，再执行 \texttt{generate --validate}；
  \item 保存 run manifest、输入、配置、中间报告、验证结果、最终 STL 和必要图片。
\end{enumerate}

\section{问题定位顺序}

\begin{description}
  \item[锚点落到导板顶部] 检查牙合轴符号、邻牙参考切向、射线角度、外壁出口法向判定和
  导板连通分量，而不是退回轨迹百分比。
  \item[U/背 U 两点在同一侧] 检查 \texttt{local\_outward}、射线符号和 FDI 左右映射，
  不要用渲染图屏幕左右替代牙弓语义。
  \item[低梁异常弯折] 检查局部下潜两侧是否继承各自的外层曲线切向，以及 P 是否过深。
  \item[导孔堵塞] 检查当前导管模式、复切顺序、cutter 轴向范围和体素尺寸。
  \item[观察窗消失] 检查最终 profile cutter 是否复切、指纹是否失效、报告 QA 是否全真。
  \item[梁被切断] 分别检查牙列保护体和手机包络，不要只看最终网格表面。
  \item[末端方向反向] 检查 FDI 方向、参考邻牙到末端牙向量和叉积符号。
\end{description}

\section{建议的正式交付清单}

\begin{itemize}
  \item 冻结的源码版本与依赖锁文件；
  \item 病例 JSON、YAML 和所有输入 STL/JSON 的摘要；
  \item tooth mapping、观察窗和手机包络的过程报告；
  \item 锚点、切口、连接梁和最终四视图；
  \item 最终 STL 的 SHA-256、拓扑统计和完整 ValidationResult；
  \item 人工设计复核记录、制造参数和实体试戴/质检记录；
  \item 本报告及其可编译 LaTeX 源工程。
\end{itemize}

\begin{takeaway}
项目交接的重点不是记住每个默认值，而是保持三条链条完整：病例语义能追溯到锚点，锚点能
追溯到最终实体，最终实体能追溯到独立验证和人工审核。
\end{takeaway}

